\documentclass[../differential_geometry.tex]{subfiles}
\begin{document}
\section{Aula 09 - de Fevereiro, 2026}
\subsection{Motivações}
\begin{itemize}
	\item Reparametrização;
	\item Superfícies Especiais;
	\item Plano Tangente.
\end{itemize}
\subsection{Superfícies Especiais}
\begin{def*}
	Seja \(f:\mathbb{R}^{3}\rightarrow \mathbb{R}\) uma função suave; um ponto p não singular em \(\mathbb{R}^{3}\) é um \textbf{ponto crítico de f} se \(\mathrm{d}f_{p}:\mathbb{R}^{3}\rightarrow \mathbb{R}\) \textit{não} é sobrejetora, ou seja, se \(\nabla f(p)=(0,0,0)\).
	Se p é um ponto crítico, diremos que \(f(p)\) é um \textbf{valor crítico de f}.

	Por outro lado, caso p não seja um ponto crítico, então diremos que p é um \textbf{ponto regular de f}. \(\square\)
\end{def*}
\begin{def*}
	Seja c um número real; então, c é um \textbf{valor regular de f} se, para todo p em \(\mathbb{R}^{3}\) tal que \(f(p)=c\), p é um valor regular. \(\square\)
\end{def*}
Equivalentemente, um valor crítico é um número real c tal que existe p ponto crítico de f com \(f(p)\) satisfazendo \(f(p)=c\).

\end{document}

